% arXiv submission: SYN-02
% Title: The Governance Impossibility Thesis: Formal Limits of
%        Self-Describing, Self-Measuring, Self-Organizing Systems
% Author: A. Padavano (Independent researcher)
% Primary category: cs.AI (Artificial Intelligence)
% Cross-list: cs.CY, cs.LO, econ.TH
% MSC 2020: 91B14, 03B25, 68T01, 91B06
% License: CC BY 4.0
% Date: 2026

\documentclass[12pt,a4paper]{article}

\usepackage[utf8]{inputenc}
\usepackage[T1]{fontenc}
\usepackage{amsmath,amssymb,amsthm}
\usepackage{hyperref}
\usepackage[sort&compress,numbers]{natbib}
\usepackage{booktabs}
\usepackage{graphicx}
\usepackage{alltt}
\usepackage{microtype}
\usepackage[margin=1in]{geometry}

\hypersetup{
  colorlinks=true,
  linkcolor=blue,
  citecolor=blue,
  urlcolor=blue
}

\newtheorem{definition}{Definition}
\newtheorem{theorem}{Theorem}
\newtheorem{proposition}{Proposition}

\title{The Governance Impossibility Thesis:\\Formal Limits of Self-Describing,\\Self-Measuring, Self-Organizing Systems}

\author{A.~Padavano\\
  \textit{Independent researcher, Studium Generale ORGANVM}}

\date{2026}

\begin{document}

\maketitle

\begin{abstract}
Self-governing systems---organizations, platforms, software architectures, and
AI-augmented institutions that attempt to describe their own structure, measure
their own performance, and organize their own coordination---confront a
convergence of impossibility results from three independent traditions. From
mathematical logic: G\"{o}del's incompleteness theorems, Tarski's
undefinability, and Rice's theorem establish that no sufficiently expressive
system can completely describe or decide its own semantic properties. From
social choice theory: Arrow's impossibility theorem and the
Gibbard-Satterthwaite theorem show that no aggregation procedure satisfies
minimal rationality constraints while faithfully reflecting diverse
preferences. From measurement theory: Goodhart's law and Campbell's law show
that metrics used for governance are corrupted by the governance they enable.
We synthesize these traditions and propose the Governance Trilemma: a
self-governing system cannot simultaneously achieve completeness (governing all
of itself), consistency (non-contradictory rules), and measurability (valid
outcome assessment). Any two can be approximated, but the third must be
sacrificed. We claim this is not a theorem in the G\"{o}delian sense but a
design constraint exhibiting the same structural signature---the impossibility
of simultaneously satisfying three desiderata within a closed system. We
derive five design responses---staged governance, hybrid topology,
translation-aware design, psychometrically calibrated measurement, and the
human-in-the-loop as incompleteness response---and apply them to a governance
architecture case study. The paper contributes a novel structural
impossibility argument for governance theory and a design framework for
systems that work within their own fundamental limits.
\end{abstract}

\noindent\textbf{Keywords:} impossibility theorems, self-governance,
G\"{o}del's incompleteness, Arrow's impossibility theorem, Goodhart's law,
construct validity, governance trilemma, mechanism design, AI governance

\bigskip

\section{Introduction}

Three intellectual traditions, developed independently across the twentieth
century, converge on a single unsettling conclusion: self-governing systems
face limits that are not engineering failures amenable to better tooling but
structural constraints rooted in the nature of self-reference, collective
decision-making, and reflexive measurement.

The first tradition is \emph{logical}. G\"{o}del~\cite{Godel1931}
demonstrated that any consistent formal system capable of encoding basic
arithmetic contains true statements it cannot prove. Tarski~\cite{Tarski1933}
showed that truth for such a system cannot be defined from within.
Rice~\cite{Rice1953} generalized these results to computation: no algorithm
can decide any non-trivial semantic property of programs.

The second tradition is \emph{social-choice-theoretic}.
Arrow~\cite{Arrow1950,Arrow1951} proved that no ranked-choice aggregation
procedure can simultaneously satisfy unrestricted domain, Pareto efficiency,
independence of irrelevant alternatives, and non-dictatorship. The
Gibbard-Satterthwaite theorem~\cite{Gibbard1973,Satterthwaite1975} extends
this to strategy-proofness.

The third tradition is \emph{measurement-theoretic}. Goodhart~\cite{Goodhart1975}
observed that ``when a measure becomes a target, it ceases to be a good
measure.'' Campbell~\cite{Campbell1979} formulated a complementary law about
the corruption of social indicators. The psychometric validity framework of
Messick~\cite{Messick1989} and Cronbach and Meehl~\cite{CronbachMeehl1955}
reveals that the constructs governance claims to measure may not correspond to
the quantities it actually captures.

This paper synthesizes these three traditions into a unified structural
argument:

\begin{quote}
\textbf{The Governance Trilemma.} A self-governing system cannot
simultaneously achieve completeness (the system can describe and govern all of
itself), consistency (the governance rules do not contradict each other), and
measurability (governance outcomes can be validly measured). Any two of these
properties can be approximated, but attaining two forces the sacrifice of the
third.
\end{quote}

We do not claim the Governance Trilemma is a theorem in the G\"{o}delian
sense---that is, a formally proven result within a precisely specified
axiomatic system. We claim it is a design constraint that exhibits the same
structural signature: the impossibility of simultaneously satisfying three
desiderata within a closed system. The trilemma is grounded in established
formal results and robust empirical regularities, and it identifies a
convergent structural pattern across these traditions.

\section{Four Sources of Impossibility}

\begin{table}[ht]
\centering
\caption{Taxonomy of impossibility sources}
\label{tab:taxonomy}
\small
\begin{tabular}{@{}llllp{3.7cm}@{}}
\toprule
Source & Level & Type & Strength & Governance Consequence \\
\midrule
Logical (G\"{o}del, Tarski, Rice) & Formal & Logical impossibility & Proven theorem & Complete self-description impossible \\
Organizational (emergence trap) & Structural & Structural impossibility & Strong empirical & Flat design $\not\Rightarrow$ flat outcomes \\
Sociomaterial (ANT, translation) & Ontological & Ontological impossibility & Philosophical argument & Meaning never perfectly transmitted \\
Measurement (Goodhart, Campbell) & Empirical & Empirical impossibility & Robust regularity & Metrics degrade under governance pressure \\
\bottomrule
\end{tabular}
\end{table}

\subsection{Logical Impossibility: The Diagonal Wall}

\textbf{G\"{o}del's first incompleteness theorem} states that any consistent,
recursively axiomatizable formal system capable of expressing basic arithmetic
is incomplete: there exist true statements in its language that it cannot
prove~\cite{Godel1931}.

\textbf{G\"{o}del's second incompleteness theorem} strengthens the first: no
consistent system satisfying these conditions can prove a formalization of its
own consistency. For governance: no governance system can certify, using its
own governance procedures, that those procedures are internally consistent.

\textbf{Tarski's undefinability theorem} establishes that truth for a
sufficiently strong formal system cannot be defined by a formula within the
system itself~\cite{Tarski1933}. Evaluation of the evaluator requires a
higher-order perspective.

\textbf{Rice's theorem} generalizes the halting problem: every non-trivial
semantic property of programs is undecidable~\cite{Rice1953}. This draws the
operational boundary for automated governance: syntactic properties can be
mechanized; semantic properties (``does this component do what it should?'')
cannot.

\textbf{L\"{o}b's theorem} constrains self-certification:
self-endorsement adds nothing~\cite{Lob1955}.

The common mechanism is the diagonal argument: when a system's descriptive
apparatus is applied to descriptions of itself, the resulting self-reference
generates objects that defeat the description.

\subsection{Organizational Impossibility: The Emergence Trap}

The \emph{compression/search tradeoff}: hierarchy functions as a compression
algorithm that reduces combinatorial complexity at the cost of information
loss; rhizomatic organization preserves connectivity at the cost of
coordination overhead.

The \emph{emergence trap}: systems designed as flat networks develop de facto
hierarchies through preferential attachment
dynamics~\cite{BarabasiAlbert1999}. This has been empirically documented
across domains (the Internet, Bitcoin, Wikipedia, open-source software). Flat
design does not guarantee flat outcomes. This resonates with Freeman's
``The Tyranny of Structurelessness''~\cite{Freeman1972}.

\subsection{Sociomaterial Impossibility: The Translation Gap}

Actor-network theory (ANT; \cite{Callon1986,Latour2005}) reveals that
translation is constitutively incomplete: governance rules do not mean the
same thing to all the entities they govern. \emph{Constitutive opacity}: AI
systems are opaque at a level that matters for governance---we cannot construct
the kind of causal account of individual decisions that governance
accountability frameworks typically presume.

\subsection{Measurement Impossibility: The Goodhart Abyss}

\textbf{Goodhart's law}~\cite{Goodhart1975,Strathern1997}: ``when a measure
becomes a target, it ceases to be a good measure.''

\textbf{Campbell's law}~\cite{Campbell1979}: the more a quantitative indicator
is used for decision-making, the more subject it is to corruption.

\textbf{Construct validity}~\cite{Messick1989,CronbachMeehl1955}: the
\emph{construct-operationalization collapse} occurs when ``quality'' is defined
as whatever the quality metrics measure. The circular definition precludes
validity assessment.

\textbf{Measurement invariance}~\cite{Meredith1993}: even metrics with
adequate construct validity in one context may lack validity in another.

An important caveat: Goodhart's law and Campbell's law do not have the same
epistemic status as the formal impossibility theorems of
Section~2.1. They are empirically observed patterns with known partial
remedies (see Section~\ref{sec:mechanism}).

\section{The Governance Trilemma}

\subsection{Completeness}

A governance system is \emph{complete} if it can describe and govern all of
its own components, processes, and behaviors. Incomplete governance invites
arbitrage.

\subsection{Consistency}

A governance system is \emph{consistent} if its rules do not contradict each
other. Inconsistent governance produces paralysis.

\subsection{Measurability}

A governance system achieves \emph{measurability} if its governance outcomes
can be validly measured---psychometrically valid, measurement-invariant, and
resistant to Goodhart corruption.

\subsection{Why You Can Only Have Two}

\textbf{Complete + Consistent $\Rightarrow$ Unmeasurable.} A comprehensive,
non-contradictory framework creates maximal conditions for Goodhart
corruption. The more complete and consistent the governance, the more
precisely actors know what to optimize for, and the more effectively they can
game the metrics.

\textbf{Complete + Measurable $\Rightarrow$ Operationally Inconsistent.}
Rice's theorem guarantees that the semantic properties the governance system
cares about are undecidable in the general case. G\"{o}del's second theorem
constrains the system's ability to prove its own consistency. The result is
not formal inconsistency but \emph{operational inconsistency}: automated
processes produce contradictory outputs at boundary cases.

\textbf{Consistent + Measurable $\Rightarrow$ Incomplete.} This is the most
practically viable combination. Consistency and measurability are achievable
for a restricted domain---but extending the framework to cover the entire
system introduces the self-referential complexity that generates
inconsistency and measurement invalidity. In practice, this is what all
successful governance systems do: they govern the domains where governance
works and leave other domains to human judgment.

\subsection{Relationship to G\"{o}del's Original Result}

The Governance Trilemma exhibits a structural pattern recognizably related to
G\"{o}del's theorems. Both concern self-referential systems. Both establish
that a desirable trio of properties cannot be jointly achieved. However, the
connection is analogical rather than deductive: we do not claim the trilemma
follows from G\"{o}del's theorems. The connection is more than metaphorical
because the same mechanism (self-application of evaluative apparatus) drives
the impossibility in both cases.

A fully formal version of the Governance Trilemma would require: (a)~a formal
definition of ``governance system'' as a mathematical object; (b)~formal
definitions of completeness, consistency, and measurability; (c)~conditions
for joint unsatisfiability; and (d)~a proof. This formalization remains future
work.

\subsection{Relationship to Arrow's Impossibility Theorem}

Arrow's theorem~\cite{Arrow1950} connects to the trilemma through the
consistency horn. The Gibbard-Satterthwaite
theorem~\cite{Gibbard1973,Satterthwaite1975} extends this to strategic
manipulation---a social-choice-theoretic analogue of Goodhart's law.

\subsection{Mechanism Design as Partial Response}
\label{sec:mechanism}

Mechanism design theory~\cite{Hurwicz1960,Maskin1999,Myerson1981} has produced
results that directly address Goodhart-type problems. The Vickrey-Clarke-Groves
(VCG) mechanism~\cite{Vickrey1961,Clarke1971,Groves1973} achieves incentive
compatibility for a significant class of allocation problems. Proper scoring
rules~\cite{Savage1971,GneitingRaftery2007} incentivize honest probability
assessments.

The trilemma's measurement horn must be stated with nuance. We claim that:
(1)~mechanism design solutions require complete specification of preference,
action, and information spaces---governance domains often resist this;
(2)~optimal mechanisms may be computationally
intractable~\cite{NisanRonen2001}; (3)~mechanism design assumes rational
agents, but governance includes non-rational actants; (4)~even if individual
metrics achieve incentive compatibility, the meta-level decision of which
metrics to use is itself subject to strategic manipulation.

\section{Design Responses to Impossibility}

\subsection{Staged Governance}

Temporal staging converts circular self-reference into sequential
cross-reference. A governance system validates its \emph{previous} state
using its \emph{current} state. Each stage is verified by a different stage.
Staging addresses G\"{o}delian impossibility and L\"{o}b's theorem directly:
the certifier is distinct from the certified.

\subsection{Hybrid Topology}

Compress where you must (hierarchy for accountability, compliance, resource
allocation). Search where you can (rhizomatic connectivity for innovation,
exploration, resilience). Make emergent hierarchy visible. Match compression
granularity to uncertainty.

\subsection{Translation-Aware Design}

Design governance artifacts as boundary objects~\cite{StarGriesemer1989}---plastic
enough to adapt to local needs, robust enough to maintain common identity.
Build feedback loops as translation-maintenance mechanisms. Acknowledge
constitutive opacity. Accept that governance is normatively underdetermined by
description.

\subsection{Psychometrically Calibrated Measurement}

Define constructs independently of their operationalizations. Test
dimensionality before compositing. Calibrate metrics using item response
theory. Test measurement invariance across contexts~\cite{Meredith1993}.
Quantify uncertainty. Design for Goodhart resistance via metric rotation and
external validation.

\textbf{Addressing the small-$N$ problem.} The statistical techniques above
presuppose sample sizes that most governance systems will not initially
achieve. Adaptations include Bayesian estimation with informative priors,
classical reliability analysis, expert panel calibration, and incremental
validation.

\subsection{The Human-in-the-Loop as Incompleteness Response}

The human provides the external perspective that G\"{o}del and Tarski demand,
the semantic judgment that Rice shows is unautomatable, the normative
authority that ANT's political vacuum leaves unfilled, and the independent
construct assessment that Goodhart corruption undermines. The human-in-the-loop
is not a concession to technological immaturity; it is a structural necessity
imposed by the impossibility results.

\subsection{The Viable System Model as Impossibility-Aware Architecture}

Beer's Viable System Model~\cite{Beer1972} maps onto the trilemma: System~3
(control) manages consistency; System~4 (intelligence) manages the
completeness boundary; System~3* (audit) manages measurability; System~5
(policy) resolves trilemma conflicts. The System~3--4 homeostat navigates
trilemma positions dynamically.

\subsection{Dynamic Trilemma Navigation}

Different lifecycle phases favor different trilemma positions:
\begin{itemize}
  \item \textbf{Early-stage:} Complete + Measurable, accepting
    inconsistency (ad hoc exceptions).
  \item \textbf{Scaling:} Consistent + Measurable, accepting
    incompleteness (the mature governance position).
  \item \textbf{Regulated:} Complete + Consistent, accepting
    unmeasurability (compliance theater).
  \item \textbf{Research:} Complete + Consistent, accepting informal
    measurement.
\end{itemize}

Migration between positions has characteristic transition costs and risks.

\section{Case Study: ORGANVM Governance Architecture}

The ORGANVM system governs approximately 117 repositories across eight organs.
It makes an implicit but consistent trilemma choice: Consistent + Measurable,
accepting Incompleteness.

\textbf{Consistency} is achieved through governance-rules.json (unidirectional
dependency flow), seed.yaml schemas, and naming conventions. \textbf{Measurability}
is achieved through composite readiness indicators and promotion state
machines. \textbf{Incompleteness} is accepted: semantic properties are
delegated to human review.

The promotion state machine (LOCAL $\to$ CANDIDATE $\to$ PUBLIC\_PROCESS $\to$
GRADUATED) implements staged governance. The IRA panel implements the Tarskian
escape---an evaluative framework richer than the governance system's formal
rules.

\textbf{Vulnerabilities:} coverage threshold gaming (Goodhart on test
coverage), seed.yaml conformance as compliance theater, omega score as
pseudo-unidimensional measure, uniform thresholds across heterogeneous
contexts.

\textbf{Recommendations:} (1)~make the trilemma choice explicit;
(2)~psychometric calibration of the scorecard; (3)~measurement invariance
testing across organs; (4)~Goodhart monitoring with metric rotation;
(5)~strengthen the IRA panel; (6)~adopt VSM recursive structure;
(7)~design for translation drift; (8)~plan for dynamic trilemma navigation.

\section{Discussion}

\subsection{The ``We Govern Fine'' Counterargument}

The trilemma does not claim existing systems are broken. It claims they are
\emph{already choosing}---implicitly selecting which trilemma vertex to
sacrifice. The contribution is making this choice explicit. The trilemma
explains the \emph{pattern} of governance failures, provides diagnostic value
for systems that do not yet exhibit failure, and generates genuinely novel
design guidance (dynamic trilemma navigation, VSM-trilemma mapping,
psychometric calibration).

\subsection{Implications for AI Governance}

AI-augmented governance intensifies every dimension of the trilemma.
Constitutively opaque actants, Goodhart compounding (AI agents producing the
measurements), and Arrovian aggregation across ontological categories all
increase the governance challenge. The human-in-the-loop faces a throughput
mismatch: AI generates governance-relevant content faster than humans can
evaluate it.

\subsection{For the Philosophy of Self-Reference}

The convergence of logical, social-choice, and measurement impossibilities
suggests that the diagonal argument has broader applicability than any single
tradition has recognized. A speculative conjecture: the impossibility results
may be aspects of a single, more general impossibility theorem.

\subsection{Formalization Roadmap}

A fully formal Governance Trilemma would require: (a)~a formal definition of
``governance system'' (perhaps a tuple $(S, R, M, E)$); (b)~formal
definitions of the three properties; (c)~conditions for joint
unsatisfiability (likely including an expressiveness threshold and a
reflexivity condition); (d)~a proof, probably constructing a governance
analogue of the G\"{o}del sentence.

\subsection{Limitations}

The case study is limited by scale (117 repositories, one operator) and
maturity. The case study is also self-referential: the system was designed
by the paper's author using principles informed by the impossibility analysis.
External application to governance systems designed without impossibility
awareness remains future work.

\section{Conclusion}

Self-governing systems confront a landscape of impossibility that is a
permanent structural feature of self-reference, collective decision-making,
and reflexive measurement. The Governance Trilemma---completeness, consistency,
and measurability cannot be jointly achieved---provides a structural framework
for understanding why governance systems persistently exhibit gaps,
contradictions, and measurement failures, and for designing systems that
acknowledge and work within these limits.

The six design responses---staged governance, hybrid topology,
translation-aware design, psychometrically calibrated measurement, the
human-in-the-loop, and dynamic trilemma navigation---are strategies for living
productively within impossibility. The art of governance, like the art of
proof, lies not in achieving completeness but in knowing where completeness
ends and judgment begins.

\section*{Acknowledgments}

This work was conducted as part of the Studium Generale ORGANVM research
program. The author acknowledges the use of AI language models (Claude,
Anthropic) as research assistants for literature synthesis, structural
drafting, and adversarial review via the Triadic Review Protocol. All
intellectual claims, arguments, and editorial decisions are the author's own.

\medskip
\noindent\textcopyright~2026 A.~Padavano. This work is licensed under a
Creative Commons Attribution 4.0 International License (CC BY 4.0).

\begin{thebibliography}{99}

\bibitem{Arrow1950}
K.~J.~Arrow.
\newblock A difficulty in the concept of social welfare.
\newblock \emph{J.\ Political Economy}, 58(4):328--346, 1950.

\bibitem{Arrow1951}
K.~J.~Arrow.
\newblock \emph{Social Choice and Individual Values}.
\newblock Yale Univ.\ Press, 1951.

\bibitem{BarabasiAlbert1999}
A.-L.~Barab\'{a}si and R.~Albert.
\newblock Emergence of scaling in random networks.
\newblock \emph{Science}, 286(5439):509--512, 1999.

\bibitem{Beer1972}
S.~Beer.
\newblock \emph{Brain of the Firm: The Managerial Cybernetics of Organization}.
\newblock Allen Lane, 1972.

\bibitem{Callon1986}
M.~Callon.
\newblock Some elements of a sociology of translation.
\newblock In J.~Law, ed., \emph{Power, Action and Belief}, pp.~196--233. Routledge, 1986.

\bibitem{Campbell1979}
D.~T.~Campbell.
\newblock Assessing the impact of planned social change.
\newblock \emph{Evaluation and Program Planning}, 2(1):67--90, 1979.

\bibitem{Clarke1971}
E.~H.~Clarke.
\newblock Multipart pricing of public goods.
\newblock \emph{Public Choice}, 11(1):17--33, 1971.

\bibitem{CronbachMeehl1955}
L.~J.~Cronbach and P.~E.~Meehl.
\newblock Construct validity in psychological tests.
\newblock \emph{Psychological Bulletin}, 52(4):281--302, 1955.

\bibitem{Freeman1972}
J.~Freeman.
\newblock The tyranny of structurelessness.
\newblock \emph{The Second Wave}, 2(1), 1972.

\bibitem{Gibbard1973}
A.~Gibbard.
\newblock Manipulation of voting schemes: A general result.
\newblock \emph{Econometrica}, 41(4):587--601, 1973.

\bibitem{GneitingRaftery2007}
T.~Gneiting and A.~E.~Raftery.
\newblock Strictly proper scoring rules, prediction, and estimation.
\newblock \emph{J.\ Amer.\ Statistical Assoc.}, 102(477):359--378, 2007.

\bibitem{Godel1931}
K.~G\"{o}del.
\newblock \"{U}ber formal unentscheidbare S\"{a}tze der Principia Mathematica und verwandter Systeme~I.
\newblock \emph{Monatshefte f\"{u}r Mathematik und Physik}, 38:173--198, 1931.

\bibitem{Goodhart1975}
C.~A.~E.~Goodhart.
\newblock Problems of monetary management: The U.K.\ experience.
\newblock In A.~S.~Courakis, ed., \emph{Inflation, Depression, and Economic Policy in the West}, pp.~116--120. Rowman \& Littlefield, 1975.

\bibitem{Groves1973}
T.~Groves.
\newblock Incentives in teams.
\newblock \emph{Econometrica}, 41(4):617--631, 1973.

\bibitem{Hurwicz1960}
L.~Hurwicz.
\newblock Optimality and informational efficiency in resource allocation processes.
\newblock In K.~J.~Arrow et al., eds., \emph{Mathematical Methods in the Social Sciences}, pp.~27--46. Stanford Univ.\ Press, 1960.

\bibitem{Latour2005}
B.~Latour.
\newblock \emph{Reassembling the Social: An Introduction to Actor-Network-Theory}.
\newblock Oxford Univ.\ Press, 2005.

\bibitem{Lob1955}
M.~H.~L\"{o}b.
\newblock Solution of a problem of Leon Henkin.
\newblock \emph{J.\ Symbolic Logic}, 20(2):115--118, 1955.

\bibitem{Maskin1999}
E.~Maskin.
\newblock Nash equilibrium and welfare optimality.
\newblock \emph{Review of Economic Studies}, 66(1):23--38, 1999.

\bibitem{Meredith1993}
W.~Meredith.
\newblock Measurement invariance, factor analysis and factorial invariance.
\newblock \emph{Psychometrika}, 58(4):525--543, 1993.

\bibitem{Messick1989}
S.~Messick.
\newblock Validity.
\newblock In R.~L.~Linn, ed., \emph{Educational Measurement}, 3rd ed., pp.~13--103. ACE/Macmillan, 1989.

\bibitem{Myerson1981}
R.~B.~Myerson.
\newblock Optimal auction design.
\newblock \emph{Mathematics of Operations Research}, 6(1):58--73, 1981.

\bibitem{NisanRonen2001}
N.~Nisan and A.~Ronen.
\newblock Algorithmic mechanism design.
\newblock \emph{Games and Economic Behavior}, 35(1--2):166--196, 2001.

\bibitem{Rice1953}
H.~G.~Rice.
\newblock Classes of recursively enumerable sets and their decision problems.
\newblock \emph{Trans.\ Amer.\ Math.\ Soc.}, 74(2):358--366, 1953.

\bibitem{Satterthwaite1975}
M.~A.~Satterthwaite.
\newblock Strategy-proofness and Arrow's conditions.
\newblock \emph{J.\ Economic Theory}, 10(2):187--217, 1975.

\bibitem{Savage1971}
L.~J.~Savage.
\newblock Elicitation of personal probabilities and expectations.
\newblock \emph{J.\ Amer.\ Statistical Assoc.}, 66(336):783--801, 1971.

\bibitem{Scott1998}
J.~C.~Scott.
\newblock \emph{Seeing Like a State}.
\newblock Yale Univ.\ Press, 1998.

\bibitem{StarGriesemer1989}
S.~L.~Star and J.~R.~Griesemer.
\newblock Institutional ecology, `translations' and boundary objects.
\newblock \emph{Social Studies of Science}, 19(3):387--420, 1989.

\bibitem{Strathern1997}
M.~Strathern.
\newblock `Improving ratings': Audit in the British university system.
\newblock \emph{European Review}, 5(3):305--321, 1997.

\bibitem{Tarski1933}
A.~Tarski.
\newblock Pojecie prawdy w jezykach nauk dedukcyjnych.
\newblock \emph{Prace Towarzystwa Naukowego Warszawskiego}, 34, 1933.

\bibitem{Vickrey1961}
W.~Vickrey.
\newblock Counterspeculation, auctions, and competitive sealed tenders.
\newblock \emph{J.\ Finance}, 16(1):8--37, 1961.

\end{thebibliography}

\end{document}
